\documentclass[a4paper,12pt]{report}

\usepackage{ucs}
\usepackage[utf8x]{inputenc} % Input encoding for Greek characters
\usepackage[greek,english]{babel} % Language support
\usepackage{siunitx}
\newcommand{\en}{\selectlanguage{english}}
\newcommand{\gr}{\selectlanguage{greek}}
% \usepackage{algorithm2e}
% \usepackage{algorithm}
% \usepackage{algorithmic}
\usepackage{enumitem}
\usepackage{float}
\usepackage{amsmath}
\usepackage{graphicx} % For including images
\usepackage{titlesec} % Custom title formatting
\usepackage{fancyhdr} % For custom headers and footers
\usepackage{geometry} % For adjusting page margins

% Adjust the page margins to make content wider
\geometry{top=2.5cm, bottom=2.5cm, left=2.5cm, right=2.5cm}

% Redefine chapter formatting to make it smaller
\titleformat{\chapter}[display]
    {\normalfont\LARGE\bfseries} % Smaller size and bold for chapter heading
    {\chaptername\ \thechapter} % Chapter number format
    {15pt} % Space between chapter number and title
    {\bfseries} % Smaller size and bold for chapter title
\begin{document}

\begin{titlepage}
    \centering
    \vspace*{-3cm}
    % University logo
    \includegraphics[width=1\textwidth]{auth_logo.png} % Replace with your actual logo file

    % University name in Greek
    \textbf{\gr ΑΡΙΣΤΟΤΕΛΕΙΟ ΠΑΝΕΠΙΣΤΗΜΙΟ ΘΕΣΣΑΛΟΝΙΚΗΣ}
    \vspace{2cm}

    % Document title and subtitle in Greek
    \LARGE\textbf{\gr Προσομοίωση και Μοντελοποίηση Δυναμικών Συστημάτων Αναφορά} \\
    \Large\normalfont{\gr Εργασία 2} \\
    \vspace{4cm}

    \gr
    \large
    \textbf{Διακολουκάς Δημήτριος} \\
    \textbf{AEM 10642}
    \vspace{2.5cm}

    \en
    \textit{Email: ddiakolou@ece.auth.gr}
\end{titlepage}

\gr
\tableofcontents

\chapter{Σύστημα μάζας–ελατηρίου–αποσβεστήρα}
\section{\gr Θέμα 1 Ανάλυση προβλήματος και εισαγωγή}
\label{chap:intro}

Στο \textbf{Θέμα 1} μελετούμε το κλασικό \emph{σύστημα μάζας–ελατηρίου–αποσβεστήρα}, το οποίο διεγείρεται από
εξωτερική δύναμη~$u(t)$.  
Η διαφορική εξίσωση που περιγράφει τη δυναμική του συστήματος είναι:

\begin{equation*}
m\,\ddot{x}(t) + b\,\dot{x}(t) + k\,x(t) = u(t),
\quad \text{με } m > 0,\; b > 0,\; k > 0
\end{equation*}

\hspace{-0.6cm}όπου
\begin{center}
\setlength{\fboxsep}{6pt}
\fbox{%
\begin{tabular}{@{}l|p{7.5cm}@{}}
\multicolumn{2}{c}{\textbf{Περιγραφή μεταβλητών}} \\[4pt]
$x(t)\,[\text{\en m\gr}]$         & η μετατόπιση της μάζας \\ \hline
$m\,[\text{\en kg\gr}]$           & η μάζα (άγνωστη) \\ \hline
$b\,[\text{\en N$\cdot$s/m\gr}]$  & ο συντελεστής απόσβεσης (άγνωστος) \\ \hline
$k\,[\text{\en N/m\gr}]$          & η σταθερά του ελατηρίου (άγνωστη) \\ \hline
$u(t)\,[\text{\en N\gr}]$         & η εξωτερική δύναμη / είσοδος ελέγχου
\end{tabular}%
}
\end{center}

\paragraph{Πειραματικές τιμές:}  
Στα αριθμητικά πειράματα θεωρούμε τις «πραγματικές»
(αλλά άγνωστες στον εκτιμητή) τιμές
\[
m_\star = 1.315,\quad
b_\star = 0.225,\quad
k_\star = 0.725 .
\]

\medskip

\subsection*{\gr Στόχοι του Θέματος}

\begin{enumerate}
  \item \textbf{Μέρος(α) — Εκτιμητής κλίσης \en (Gradient estimator\gr).}
        \begin{itemize}
          \item Σχεδιάζουμε εκτιμητή πραγματικού χρόνου για τα~$m,b,k$
                χρησιμοποιώντας τη μέθοδο κλίσης.
          \item Εξετάζουμε δύο διεγέρσεις:
                (\en i\gr)~σταθερό σήμα $u(t)=2.5$ και
                (\en ii\gr)~ημιτονοειδές σήμα $u(t)=2.5\sin t$.
          \item Παρουσιάζουμε γραφήματα των $x(t)$, $\hat x(t)$,
                του σφάλματος~$e_x(t)=x(t)-\hat x(t)$
                και των εκτιμήσεων $\hat m(t),\hat b(t),\hat k(t)$.
        \end{itemize}

  \item \textbf{Μέρος (β) — Εκτιμητής \en Lyapunov \gr \emph{παράλληλης} και
        \emph{μικτής} \en (series–parallel\gr) δομής.}
        \begin{itemize}
          \item Υλοποιούμε εκτιμητή βασισμένο σε αναλυτική κατασκευή
                συνάρτησης \en Lyapunov\gr, για $u(t)=2.5\sin t$.
          \item Συγκρίνουμε παράλληλη \en vs\gr.\ μικτή δομή ως προς
                ταχύτητα σύγκλισης και
                ακρίβεια των $\hat m,\hat b,\hat k$.
        \end{itemize}

  \item \textbf{Μέρος (γ) — Επίδραση μετρούμενου θορύβου.}
        \begin{itemize}
          \item Προσθέτουμε θόρυβο μέτρησης
                $h(t)=h_0\sin(2\pi f_0 t)$ στην έξοδο $x(t)$
                (με $h_0=0.25$ και $f_0=20$\,\en Hz\gr).
          \item Επαναλαμβάνουμε το Μέρος (β) και
                μελετούμε πώς το πλάτος θορύβου
                επηρεάζει το σφάλμα εκτίμησης.
        \end{itemize}
\end{enumerate}

\medskip

\subsection*{\gr Συντομογραφίες και υποθέσεις}

\begin{itemize}
  \item Τα σήματα $x(t)$ και $u(t)$ θεωρούνται διαθέσιμα
        (μετρήσιμα σε πραγματικό χρόνο).
  \item Η λύση υλοποιείται σε κώδικα στο \en \textsc{MATLAB} \gr
        με βήμα $h=10^{-2}\,\mathrm{s}$ για διάστημα $20$\,\en s\gr,
        ώστε να καταγραφεί η δυναμική σύγκλισης.
  \item Όπου απαιτείται, χρησιμοποιείται η τυπική μορφή
        \(
        \dot z = Az + Bu,\;
        y = Cz
        \)
        για ανάλυση σταθερότητας του σφάλματος εκτίμησης.
\end{itemize}

\medskip
\hspace{-0.6cm }Με την παραπάνω εισαγωγή διασαφηνίζονται τα φυσικά μεγέθη,
οι στόχοι προσομοίωσης και οι επιμέρους απαιτήσεις
για τα Μέρη (α)–(β)–(γ) του Θέματος 1.
Στα επόμενα \en sections \gr παρουσιάζονται
\textbf{\en (i\gr)}~οι υλοποιήσεις σε κώδικα \en \textsc{MATLAB}\gr,
\textbf{\en (ii\gr)}~τα αποτελέσματα προσομοίωσης
και \textbf{\en (iii\gr)}~η συζήτηση / ερμηνεία
των γραφημάτων σύγκλισης των εκτιμητών.

\bigskip

\section{Θέμα 1(α) - Σχεδίαση εκτιμητή με μέθοδο κλίσης}

Για την επίλυση του προβλήματος πραγματικού χρόνου εκτίμησης των άγνωστων παραμέτρων 
$m$, $b$ και $k$ του συστήματος, εφαρμόζεται η μέθοδος \emph{κλίσης (\en gradient method\gr)} 
σε κατάλληλη φιλτραρισμένη μορφή. Το αρχικό δυναμικό μοντέλο είναι:

\begin{equation*}
m\,\ddot{x}(t) + b\,\dot{x}(t) + k\,x(t) = u(t),
\quad \text{με } m>0,\; b>0,\; k>0
\end{equation*}

\hspace{-0.6cm}Η μέθοδος υλοποιείται με βάση την ιδέα ότι η δυναμική της επιτάχυνσης μπορεί να γραφεί ως:

\[
\ddot{x}(t) = \theta_1\,\dot{x}(t) + \theta_2\,x(t) + \theta_3\,u(t)
\]

\hspace{-0.6cm}όπου
\[
\theta_1 = -\frac{b}{m},\quad
\theta_2 = -\frac{k}{m},\quad
\theta_3 = \frac{1}{m}
\]

\medskip

\subsection*{Αντικειμενικός σκοπός:}
Ο στόχος είναι να εκτιμηθούν οι $\theta_1$, $\theta_2$, $\theta_3$ μέσω φίλτρων δεύτερης τάξης 
και του σφάλματος εκτίμησης $e(t) = x(t) - \hat{x}(t)$, όπου η εκτίμηση $\hat{x}(t)$ δίνεται ως:

\[
\hat{x}(t) = \theta_1(t)\,j_1(t) + \theta_2(t)\,j_2(t) + \theta_3(t)\,j_3(t)
\]

\hspace{-0.6cm}Οι $j_i(t)$ αποτελούν φιλτραρισμένες εκδοχές των $\dot{x}(t)$, $x(t)$ και $u(t)$ αντίστοιχα, 
οι οποίες ικανοποιούν διαφορικές εξισώσεις τύπου:

\[
\ddot{j}_i(t) + (p_1 + p_2)\,\dot{j}_i(t) + p_1p_2\,j_i(t) = \text{σήμα εισόδου}
\]

\hspace{-0.6cm}όπου $p_1$, $p_2$ είναι προκαθορισμένοι πόλοι φίλτρου (\en filter poles\gr) του συστήματος.

\medskip

\subsection*{Νόμος προσαρμογής \en (gradient update law\gr):}
Για κάθε $\theta_i$ εφαρμόζεται νόμος εκμάθησης της μορφής:

\[
\dot{\theta}_i(t) = \gamma_i\,e(t)\,j_i(t),
\quad i=1,2,3
\]

\hspace{-0.6cm}όπου $\gamma_i$ είναι θετικές σταθερές προσαρμογής (\en learning rates\gr). 
Αυτή η προσαρμοστική διαδικασία επιτρέπει τη σύγκλιση των $\theta_i(t)$ προς τις 
πραγματικές τιμές μέσω της ελαχιστοποίησης του σφάλματος $e(t)$.

\medskip

\subsection*{Ανάκτηση φυσικών παραμέτρων:}
Αφού εκτιμηθούν οι παράμετροι $\theta_i(t)$, οι πραγματικές παράμετροι του συστήματος ανακτώνται ως:

\[
\hat{m}(t) = \frac{1}{\theta_3(t)},
\quad
\hat{b}(t) = -\theta_1(t) \cdot \hat{m}(t),
\quad
\hat{k}(t) = -\theta_2(t) \cdot \hat{m}(t)
\]

\medskip

\subsection*{Διερεύνηση δύο περιπτώσεων εισόδου:}
Η διαδικασία εφαρμόζεται για δύο μορφές εισόδου:
\begin{enumerate}
  \item \textbf{Σταθερή είσοδος:} $u(t) = 2.5$
  \item \textbf{Ημιτονοειδής είσοδος:} $u(t) = 2.5\sin(t)$
\end{enumerate}

\hspace{-0.6cm}Και στις δύο περιπτώσεις εκτελείται χρονική ολοκλήρωση $20\,\text{\en s \gr}$ με 
βήμα $0.01\,\text{\en s \gr}$ και συγκρίνονται τα διαγράμματα της πραγματικής $x(t)$ με την εκτιμώμενη $\hat{x}(t)$, 
καθώς και οι τελικές εκτιμήσεις για $m$, $b$, $k$.

\medskip

\subsection*{Αποτελέσματα Eκτίμησης και \en Plots \gr:}
\begin{center}
\begin{tabular}{@{}l l@{\hskip 1em}l@{}}
\multicolumn{3}{l}{\textbf{\hspace{-1cm} - Περίπτωση:} $u(t) = 2.5$} \\
Μάζα:               & \en true = 1.3150,     & \en estimated = 1.3129 \gr \\
Απόσβεση:           & \en true = 0.2250,     & \en estimated = 0.2359 \gr \\
Σταθερά ελατηρίου:  & \en true = 0.7250,     & \en estimated = 0.7246 \gr \\
\\
\multicolumn{3}{l}{\textbf{\hspace{-1cm} - Περίπτωση:} $u(t) = 2.5\sin(t)$} \\
Μάζα:               & \en true = 1.3150,     & \en estimated = 1.3114 \gr \\
Απόσβεση:           & \en true = 0.2250,     & \en estimated = 0.2246 \gr \\
Σταθερά ελατηρίου:  & \en true = 0.7250,     & \en estimated = 0.7216 \gr \\
\end{tabular}
\end{center}

\begin{figure}[H]
  \centering
  \includegraphics[width=0.85\textwidth]{latex_and_pdf/1.a.1.png}
  \caption{Πραγματική και εκτιμώμενη μετατόπιση $x(t)$ και $\hat{x}(t)$ για σταθερή είσοδο}
\end{figure}

\begin{figure}[H]
  \centering
  \includegraphics[width=0.85\textwidth]{latex_and_pdf/1.a.2.png}
  \caption{Σφάλμα εκτίμησης $e(t)$ για σταθερή είσοδο $u(t)=2.5$}
\end{figure}

\begin{figure}[H]
  \centering
  \includegraphics[width=0.85\textwidth]{latex_and_pdf/1.a.3.png}
  \caption{Εκτιμήσεις παραμέτρων $m(t)$, $b(t)$, $k(t)$ για σταθερή είσοδο}
\end{figure}

\begin{figure}[H]
  \centering
  \includegraphics[width=0.85\textwidth]{latex_and_pdf/1.a.4.png}
  \caption{Πραγματική και εκτιμώμενη μετατόπιση $x(t)$ και $\hat{x}(t)$ για ημιτονοειδή είσοδο}
\end{figure}

\begin{figure}[H]
  \centering
  \includegraphics[width=0.85\textwidth]{latex_and_pdf/1.a.5.png}
  \caption{Σφάλμα εκτίμησης $e(t)$ για ημιτονοειδή είσοδο $u(t)=2.5\sin(t)$}
\end{figure}

\begin{figure}[H]
  \centering
  \includegraphics[width=0.85\textwidth]{latex_and_pdf/1.a.6.png}
  \caption{Εκτιμήσεις παραμέτρων $m(t)$, $b(t)$, $k(t)$ για ημιτονοειδή είσοδο}
\end{figure}

\medskip

\subsection*{Παρατηρήσεις}
\begin{itemize}
\item Οι εκτιμήσεις των παραμέτρων ($\hat{m}$, $\hat{b}$, $\hat{k}$) συγκλίνουν γρήγορα  (\en 20 sec\gr) στις πραγματικές τιμές και για τις δύο περιπτώσεις εισόδου με κατάλληλη επιλογή φίλτρου και σωστό \en tuning \gr των παραμέτρων γ.

\item Το σφάλμα εκτίμησης $e_x(t) = x(t) - \hat{x}(t)$ μειώνεται γρήγορα σε σχεδόν μηδενικά επίπεδα, δείχνοντας καλή συμπεριφορά του εκτιμητή.

\item Η χρήση φίλτρου δεύτερης τάξης με πόλους στο $s^2 - s + 1$ εξασφαλίζει επαρκή απόσβεση και μηδενισμό των παραγώγων, προσφέροντας ομαλές εκτιμήσεις ακόμα και σε μη ιδανικές συνθήκες.

\end{itemize}

\hspace{-0.6cm}Ο αντίστοιχος κώδικας \en MATLAB \gr βρίσκεται στο αρχείο \en Part1a.m\gr.

\bigskip

\section{Θέμα 1(β) - Εκτιμητής \en Lyapunov \gr (παράλληλης και μεικτής δομής)}

Για την επίλυση του προβλήματος εκτίμησης των αγνώστων παραμέτρων $m$, $b$, $k$ 
του συστήματος, εφαρμόζεται η μέθοδος \en Lyapunov \gr τόσο σε \textbf{παράλληλη δομή} όσο και σε 
\textbf{μεικτή (\en series-parallel\gr)}. Η είσοδος είναι ημιτονοειδής:
\[
u(t) = 2.5 \sin(t)
\]

\medskip

\subsection*{Παράλληλη Δομή:}

Η παράλληλη δομή βασίζεται σε εκτίμηση του συστήματος της μορφής:
\[
\hat{x}_1(t) = x_2(t),\qquad
\hat{x}_2(t) = \hat{a}_{21}(t) x_1(t) + \hat{a}_{22}(t) x_2(t) + \hat{b}_2(t) u(t)
\]

\hspace{-0.6cm}και χρησιμοποιεί το σφάλμα εκτίμησης:
\[
e_2(t) = x_2(t) - \hat{x}_2(t)
\]

\hspace{-0.6cm}Ο νόμος προσαρμογής είναι της μορφής:
\[
\dot{\hat{a}}_{21} = \gamma_1 e_2 x_1, \qquad
\dot{\hat{a}}_{22} = \gamma_2 e_2 x_2, \qquad
\dot{\hat{b}}_2    = \gamma_3 e_2 u
\]

\hspace{-0.6cm}Οι εκτιμήσεις των φυσικών παραμέτρων προκύπτουν ως:
\[
\hat{m} = \frac{1}{\hat{b}_2}, \qquad
\hat{b} = -\hat{a}_{22} \cdot \hat{m}, \qquad
\hat{k} = -\hat{a}_{21} \cdot \hat{m}
\]

\medskip

\subsection*{Μεικτή Δομή \en (Series-Parallel\gr):}

Η μεικτή δομή βελτιώνει τη δυναμική παρατήρηση εισάγοντας όρους διόρθωσης στο σύστημα εκτίμησης:
\[
\begin{aligned}
\dot{x}_1 &= x_2 \\
\dot{x}_2 &= a_{21} x_1 + a_{22} x_2 + b_2 u \\
\dot{\hat{x}}_1 &= x_2 + \theta_{m1} e_1 + \theta_{m2} e_2 \\
\dot{\hat{x}}_2 &= \hat{a}_{21} x_1 + \hat{a}_{22} x_2 + \hat{b}_2 u + \theta_{m3} e_1 + \theta_{m4} e_2
\end{aligned}
\]

\hspace{-0.6cm}Με $e_1 = x_1 - \hat{x}_1$ και $e_2 = x_2 - \hat{x}_2$.

\medskip

\subsection*{Τιμές προσαρμογής και διόρθωσης:}

Οι σταθερές προσαρμογής (\en learning rates\gr) και οι όροι διόρθωσης έχουν επιλεχθεί ως εξής:

\begin{itemize}
  \item \textbf{Παράλληλη δομή:}
  \begin{itemize}
    \item $\gamma_1 = 0.0475$
    \item $\gamma_2 = 0.0454$
    \item $\gamma_3 = 0.0167$
  \end{itemize}

  \item \textbf{Μεικτή δομή \en (Series–Parallel\gr):}
  \begin{itemize}
    \item $\gamma_1 = 0.062$
    \item $\gamma_2 = 0.0597$
    \item $\gamma_3 = 1.503$
    \item Όροι διόρθωσης: $\theta_{m1} = \theta_{m2} = \theta_{m3} = \theta_{m4} = 0.2$
  \end{itemize}
\end{itemize}

\hspace{-0.6cm}Οι παραπάνω τιμές επηρεάζουν τον ρυθμό προσαρμογής των εκτιμήσεων και ρυθμίζουν την ταχύτητα σύγκλισης του εκτιμητή.

\medskip

\subsection*{Αποτελέσματα Εκτίμησης και \en Plots \gr:}

\begin{center}
\begin{tabular}{@{}l l@{\hskip 1em}l@{}}
\multicolumn{3}{l}{\textbf{\hspace{-1cm} - Παράλληλη Δομή:}} \\
Μάζα:               & \en true = 1.3150,     & \en estimated = 1.3152 \gr \\
Απόσβεση:           & \en true = 0.2250,     & \en estimated = 0.2253 \gr \\
Σταθερά ελατηρίου:  & \en true = 0.7250,     & \en estimated = 0.7076 \gr \\
\\
\multicolumn{3}{l}{\textbf{\hspace{-1cm} - Μεικτή Δομή:}} \\
Μάζα:               & \en true = 1.3150,     & \en estimated = 1.3212 \gr \\
Απόσβεση:           & \en true = 0.2250,     & \en estimated = 0.2234 \gr \\
Σταθερά ελατηρίου:  & \en true = 0.7250,     & \en estimated = 0.7175 \gr \\
\end{tabular}
\end{center}

\begin{figure}[H]
  \centering
  \includegraphics[width=0.85\textwidth]{latex_and_pdf/1.b.1.png}
  \caption{Παράλληλη δομή: $x_1(t)$ και $\hat{x}_1(t)$}
\end{figure}

\begin{figure}[H]
  \centering
  \includegraphics[width=0.85\textwidth]{latex_and_pdf/1.b.2.png}
  \caption{Παράλληλη δομή: Σφάλμα εκτίμησης $e_1(t)$}
\end{figure}

\begin{figure}[H]
  \centering
  \includegraphics[width=0.85\textwidth]{latex_and_pdf/1.b.3.png}
  \caption{Παράλληλη δομή: Εκτιμήσεις $m(t)$, $b(t)$, $k(t)$}
\end{figure}

\begin{figure}[H]
  \centering
  \includegraphics[width=0.85\textwidth]{latex_and_pdf/1.b.4.png}
  \caption{Μεικτή δομή: $x_1(t)$ και $\hat{x}_1(t)$}
\end{figure}

\begin{figure}[H]
  \centering
  \includegraphics[width=0.85\textwidth]{latex_and_pdf/1.b.5.png}
  \caption{Μεικτή δομή: Σφάλμα εκτίμησης $e_1(t)$}
\end{figure}

\begin{figure}[H]
  \centering
  \includegraphics[width=0.85\textwidth]{latex_and_pdf/1.b.6.png}
  \caption{Μεικτή δομή: Εκτιμήσεις $m(t)$, $b(t)$, $k(t)$}
\end{figure}

\subsection*{Παρατηρήσεις}
\begin{itemize}
  \item Και στις δύο δομές παρατηρείται πολύ καλή προσέγγιση των πραγματικών παραμέτρων, με μικρό σφάλμα εκτίμησης.
  \item Η μεικτή δομή δίνει ελαφρώς πιο ομαλή σύγκλιση και καλύτερη δυναμική, λόγω του όρου διόρθωσης $e_1, e_2$.
  \item Οι αποκλίσεις στα πρώτα 0.5 \en sec \gr είναι φυσιολογικές και μειώνονται γρήγορα ενώ με κατάλληλο \en tuning \gr των παραμέτρων γ μπορούμε να πετύχουμε όλο και καλύτερη σύγκλιση στον χρόνο και εκτιμήσεις.
\end{itemize}

\hspace{-0.6cm}Ο αντίστοιχος κώδικας \en MATLAB \gr βρίσκεται στο αρχείο \en Part1b.m\gr.

\bigskip

\section{Θέμα 1(γ) - Εκτίμηση με παρουσία θορύβου μέτρησης}

Το ερώτημα αυτό είναι εμμέσως συνδεδεμένο με την προηγούμενη υλοποίηση στο ερώτημα 1(β). Σε αυτό το ερώτημα μελετάται η συμπεριφορά των εκτιμητών \en Lyapunov \gr (παράλληλης και μεικτής δομής) όταν η μετρούμενη έξοδος $x(t)$ περιέχει πρόσθετο περιοδικό θόρυβο μέτρησης της μορφής:
\[
\eta(t) = \eta_0 \sin(2\pi f_0 t)
\]
με $\eta_0 = 0.25$ και $f_0 = 20\,\text{\en Hz\gr}$. Το σήμα θορύβου προστίθεται στην $x_1(t)$ όπως φαίνεται από τις συναρτήσεις υλοποίησης του εκτιμητή.

\medskip

\subsection*{Παράλληλη Δομή (με θόρυβο)}

Η δυναμική είναι όπως και στο ερώτημα (β), με προσθήκη του θορύβου στην είσοδο της εξίσωσης σφάλματος:
\[
e_1(t) = (x_1(t) + \eta(t)) - \hat{x}_1(t), \qquad
e_2(t) = x_2(t) - \hat{x}_2(t)
\]

\hspace{-0.6cm}Ο νόμος προσαρμογής παραμένει:
\[
\dot{\hat{a}}_{21} = \gamma_1 e_2 x_1, \qquad
\dot{\hat{a}}_{22} = \gamma_2 e_2 x_2, \qquad
\dot{\hat{b}}_2    = \gamma_3 e_2 u
\]

\medskip

\subsection*{Μεικτή Δομή \en (Series-Parallel\gr) με θόρυβο}

Η διαφορά εδώ είναι ότι και στη $\dot{\hat{x}}_2$ εισέρχεται το σήμα $x_1 + \eta(t)$:
\[
\dot{\hat{x}}_2 = \hat{a}_{21} (x_1+\eta) + \hat{a}_{22} x_2 + \hat{b}_2 u + \theta_{m3} e_1 + \theta_{m4} e_2
\]

\hspace{-0.6cm}Οι υπόλοιποι όροι και προσαρμογές διατηρούνται.

\medskip

\subsection*{Τιμές προσαρμογής και διόρθωσης:}

\begin{itemize}
  \item \textbf{Παράλληλη δομή (με θόρυβο):}
  \begin{itemize}
    \item $\gamma_1 = 0.059$
    \item $\gamma_2 = 0.064$
    \item $\gamma_3 = 0.0191$
  \end{itemize}

  \item \textbf{Μεικτή δομή (με θόρυβο):}
  \begin{itemize}
    \item $\gamma_1 = 0.060$
    \item $\gamma_2 = 0.057$
    \item $\gamma_3 = 1.5$
    \item $\theta_{mj} = 0.2 \quad \text{για } j=1,2,3,4$
  \end{itemize}
\end{itemize}

\medskip

\subsection*{Αποτελέσματα Εκτίμησης με Θόρυβο:}

\begin{center}
\begin{tabular}{@{}l l@{\hskip 1em}l@{}}
\multicolumn{3}{l}{\textbf{\hspace{-1cm} - Παράλληλη Δομή (θόρυβος):}} \\
Μάζα:               & \en true = 1.3150,     & \en estimated = 1.3264 \gr \\
Απόσβεση:           & \en true = 0.2250,     & \en estimated = 0.2226 \gr \\
Σταθερά ελατηρίου:  & \en true = 0.7250,     & \en estimated = 0.7264 \gr \\
\\
\multicolumn{3}{l}{\textbf{\hspace{-1cm} - Μεικτή Δομή (θόρυβος):}} \\
Μάζα:               & \en true = 1.3150,     & \en estimated = 1.3190 \gr \\
Απόσβεση:           & \en true = 0.2250,     & \en estimated = 0.2083 \gr \\
Σταθερά ελατηρίου:  & \en true = 0.7250,     & \en estimated = 0.7228 \gr \\
\end{tabular}
\end{center}

\begin{figure}[H]
  \centering
  \includegraphics[width=0.85\textwidth]{latex_and_pdf/1.c.1.png}
  \caption{Παράλληλη δομή: $x_1(t)$ και $\hat{x}_1(t)$ με προσθήκη θορύβου}
\end{figure}

\begin{figure}[H]
  \centering
  \includegraphics[width=0.85\textwidth]{latex_and_pdf/1.c.2.png}
  \caption{Παράλληλη δομή: Σφάλμα εκτίμησης $e_1(t)$ (θόρυβος)}
\end{figure}

\begin{figure}[H]
  \centering
  \includegraphics[width=0.85\textwidth]{latex_and_pdf/1.c.3.png}
  \caption{Παράλληλη δομή: Εκτιμήσεις $m(t)$, $b(t)$, $k(t)$ (θόρυβος)}
\end{figure}

\begin{figure}[H]
  \centering
  \includegraphics[width=0.85\textwidth]{latex_and_pdf/1.c.4.png}
  \caption{Μεικτή δομή: $x_1(t)$ και $\hat{x}_1(t)$ με προσθήκη θορύβου}
\end{figure}

\begin{figure}[H]
  \centering
  \includegraphics[width=0.85\textwidth]{latex_and_pdf/1.c.5.png}
  \caption{Μεικτή δομή: Σφάλμα εκτίμησης $e_1(t)$ (θόρυβος)}
\end{figure}

\begin{figure}[H]
  \centering
  \includegraphics[width=0.85\textwidth]{latex_and_pdf/1.c.6.png}
  \caption{Μεικτή δομή: Εκτιμήσεις $m(t)$, $b(t)$, $k(t)$ (θόρυβος)}
\end{figure}

\medskip

\subsection*{Παρατηρήσεις}
\begin{itemize}
  \item Η παρουσία περιοδικού θορύβου μειώνει την ακρίβεια των εκτιμητών, ειδικά στην αρχική φάση.
  \item Η μεικτή δομή επιτυγχάνει ομαλότερες εκτιμήσεις έναντι της παράλληλης λόγω των διορθώσεων $e_1$, $e_2$.
  \item Οι εκτιμήσεις συγκλίνουν σε ικανοποιητικά επίπεδα παρά τον θόρυβο, με απόκλιση αμελητέα μικρή με σωστό \en tuning \gr παραμέτρων προσαρμογής και διόρθωσης.
\end{itemize}

\hspace{-0.6cm}Ο αντίστοιχος κώδικας \en MATLAB \gr βρίσκεται στο αρχείο \en Part1c.m\gr. 


\chapter{Σύστημα ελέγχου γωνίας κλίσης αεροσκάφους}
\section{\gr Θέμα 2 – Ανάλυση προβλήματος και εισαγωγή}

Στο \textbf{Θέμα 2} μελετάται το \emph{μη-γραμμικό σύστημα ελέγχου της γωνίας κλίσης (\en roll angle\gr)} ενός αεροσκάφους, με ροπή εισόδου~$u(t)$ και εξωτερικές διαταραχές~$d(t)$.  
Η διαφορική εξίσωση που περιγράφει το σύστημα είναι:

\begin{equation*}
\ddot{r}(t) = -a_1\dot{r}(t) - a_2 \sin(r(t)) + a_3\dot{r}^2(t)\sin(2r(t)) + b\,u(t) + d(t)
\end{equation*}

\hspace{-0.6cm}όπου:

\begin{center}
\setlength{\fboxsep}{6pt}
\fbox{%
\begin{tabular}{@{}l|p{7.5cm}@{}}
\multicolumn{2}{c}{\textbf{Περιγραφή μεταβλητών}} \\[4pt]
$r(t)\,[\text{\en rad\gr}]$       & η γωνία κλίσης του αεροσκάφους \\ \hline
$a_1,\,a_2,\,a_3 > 0$             & άγνωστοι συντελεστές δυναμικής \\ \hline
$b > 0$                           & άγνωστος συντελεστής ροπής ελέγχου \\ \hline
$u(t)\,[\text{\en Nm\gr}]$        & ροπή εισόδου \\ \hline
$d(t)\,[\text{\en Nm\gr}]$        & εξωτερική διαταραχή \\ \hline
$\dot{r}(t)\,[\text{\en rad/s\gr}]$ & ρυθμός μεταβολής γωνίας
\end{tabular}%
}
\end{center}

\paragraph{Πειραματικές τιμές:}

Θεωρούμε ότι οι πραγματικές τιμές (άγνωστες στον εκτιμητή) είναι:
\[
a_1 = 1.315, \quad
a_2 = 0.725, \quad
a_3 = 0.225, \quad
b = 1.175.
\]

\medskip

\subsection*{\gr Στόχοι του Θέματος}

\begin{enumerate}
  \item \textbf{Μέρος (α) — Ελεγκτής ανάδρασης με στοχοθέτηση τροχιάς.}
        \begin{itemize}
          \item Υλοποιείται ελεγκτής $u(t)$ ώστε η $r(t)$ να ακολουθεί την επιθυμητή τροχιά
                $r_d(t)$, με $r_d(0) = 0$ και $r_d(t) \to \frac{\pi}{10}$, $t \to 10$.
          \item Παρουσιάζονται τα διαγράμματα $r(t)$ και $r_d(t)$
                υπό την επίδραση του ελεγκτή.
        \end{itemize}

  \item \textbf{Μέρος (β) — Εκτίμηση παραμέτρων με χρήση της μεθόδου \en Lyapunov\gr.}
        \begin{itemize}
          \item Θεωρούμε διαθέσιμα τα σήματα $r(t)$, $\dot{r}(t)$ και $u(t)$.
          \item Υλοποιείται εκτιμητής για τις άγνωστες παραμέτρους $a_1, a_2, a_3, b$.
          \item Παρουσιάζεται η εκτίμηση $\hat{r}(t)$ και οι εκτιμήσεις $\hat{a}_1, \hat{a}_2, \hat{a}_3, \hat{b}$.
        \end{itemize}

  \item \textbf{Μέρος (γ) — Εξωτερικές διαταραχές.}
        \begin{itemize}
          \item Εισάγεται εξωτερική διαταραχή της μορφής $d(t) = 0.15\sin(0.5t)$.
          \item Επαναλαμβάνεται η διαδικασία του (β) και μελετάται η επίδραση της διαταραχής
                στις τελικές εκτιμήσεις.
        \end{itemize}
\end{enumerate}

\medskip

\subsection*{\gr Παραδοχές και προϋποθέσεις}

\begin{itemize}
  \item Οι καταστάσεις $r(t)$ και $\dot{r}(t)$ καθώς και το $u(t)$ είναι διαθέσιμες.
  \item Ο θόρυβος μέτρησης θεωρείται αμελητέος στα Μέρη (α)–(β), ενώ λαμβάνεται υπόψη στο (γ).
  \item Όλες οι εξισώσεις προσομοιώνονται σε περιβάλλον \en MATLAB\gr\ με χρονικό βήμα $dt = 10^{-4}$ και διάρκεια $20\,\text{\en s\gr}$.
\end{itemize}

\medskip
\hspace{-0.6cm }Η παραπάνω εισαγωγή συνοψίζει το πρόβλημα ελέγχου και εκτίμησης
που αντιμετωπίζεται στο Θέμα~2, καθώς και τις βασικές προδιαγραφές
προσομοίωσης και υλοποίησης. Στα επόμενα \en sections \gr παρουσιάζονται αναλυτικά
\textbf{\en (i\gr)} η υλοποίηση του ελεγκτή και του εκτιμητή, 
\textbf{\en (ii\gr)} τα αποτελέσματα, 
και \textbf{\en (iii\gr)} οι σχετικές παρατηρήσεις και συγκρίσεις.

\bigskip

\section{Θέμα 2(α) – Ελεγκτής ανάδρασης για τροχιά $r_d(t)$}

Στο \textbf{Θέμα 2(α)} μελετάται η παρακολούθηση επιθυμητής γωνιακής τροχιάς $r_d(t)$ από τη \en roll angle \gr $r(t)$ ενός αεροσκάφους, το οποίο περιγράφεται από τη μη-γραμμική δυναμική:

\[
\ddot{r}(t) = -a_1 \dot{r}(t) - a_2 \sin(r(t)) + a_3 \dot{r}^2(t) \sin(2r(t)) + b\,u(t)
\]

όπου:
\begin{itemize}
  \item $r(t)$: η γωνία κλίσης \en (roll angle) \gr σε ακτίνια,
  \item $u(t)$: η είσοδος ελέγχου (ροπή),
  \item $a_1$, $a_2$, $a_3$, $b$: θετικοί, άγνωστοι παράμετροι.
\end{itemize}

\medskip

\subsection*{Σκοπός ελέγχου}
Να σχεδιαστεί ένας μη-γραμμικός ελεγκτής $u(t)$ έτσι ώστε η έξοδος $r(t)$:
\begin{itemize}
  \item να ξεκινά από $r(0) = 0$,
  \item να φτάνει την επιθυμητή τιμή $r_{\text{\en peak\gr}} = \frac{\pi}{10} \approx 0.314\,\text{\en rad\gr}$ στο $t \approx 10$\,\text{\en s\gr},
  \item και να επιστρέφει ομαλά στο μηδέν έως $t = 20$\,\text{\en s\gr}.
\end{itemize}

\paragraph{Χρησιμοποιείται η $C^2$ συνεχής τροχιά αναφοράς:}
\[
r_d(t) = r_{\text{\en peak\gr}} \cdot \sin^2\left( \frac{\pi t}{T} \right),
\quad \text{με} \quad r_{\text{\en peak\gr}} = \frac{\pi}{10},\quad T = 20\,\text{\en s\gr}
\]

\medskip

\subsection*{Μηχανισμός ελέγχου (\en funnel control\gr)}

Για να περιορίσουμε το σφάλμα $r(t) - r_d(t)$ εντός φθίνοντος περιθωρίου, χρησιμοποιείται μια \en \textbf{parametric funnel} \gr της μορφής:
\[
\phi(t) = (\phi_0 - \phi_\infty)\,e^{-\lambda t} + \phi_\infty
\]

και ορίζονται οι μεταβλητές ελέγχου:
\[
z_1(t) = \frac{r(t) - r_d(t)}{\phi(t)},
\quad
a(t) = -k_1 \ln \left( \frac{1 + z_1(t)}{1 - z_1(t)} \right)
\]
\[
z_2(t) = \frac{\dot{r}(t) - a(t)}{\rho},
\quad
u(t) = -k_2 \ln \left( \frac{1 + z_2(t)}{1 - z_2(t)} \right)
\]

\paragraph{Επιλεγμένες παράμετροι ελεγκτή:}
\[
\phi_0 = 0.5,\quad \phi_\infty = 0.01,\quad \lambda = 0.5,\quad
\rho = 1,\quad k_1 = 2,\quad k_2 = 2
\]

Οι παραπάνω ρυθμίσεις επιτρέπουν:
\begin{itemize}
  \item \textbf{Ακριβή και ασφαλή παρακολούθηση}, χωρίς υπερβάσεις,
  \item \textbf{Γραμμική σύγκλιση} του σφάλματος προς το μηδέν,
  \item \textbf{Φυσικά φραγμένες εντολές ελέγχου} $u(t)$ χωρίς κορεσμό.
\end{itemize}

\medskip

\subsection*{Αποτελέσματα Προσομοίωσης}

Η προσομοίωση εκτελέστηκε με χρονικό βήμα $h=10^{-4}\,\text{\en s \gr}$ για $T=20\,\text{\en s\gr}$, με αρχικές συνθήκες $r(0)=0$, $\dot{r}(0)=0$.

\begin{figure}[H]
  \centering
  \includegraphics[width=0.85\textwidth]{latex_and_pdf/2.a.1.png}
  \caption{Γωνία ρόλ $r(t)$ και επιθυμητή τροχιά $r_d(t)$ — απόκριση κλειστού βρόχου}
\end{figure}

\begin{figure}[H]
  \centering
  \includegraphics[width=0.85\textwidth]{latex_and_pdf/2.a.2.png}
  \caption{Σήμα ελέγχου $u(t)$ που παράγεται από τον μη-γραμμικό ελεγκτή}
\end{figure}

\paragraph{Αριθμητικό αποτέλεσμα:}
\en
\[
\text{Reach } r_{\text{peak}} = 0.314\,\text{rad} \text{ and return to } 0 \text{ in } 20\,\text{s}.
\]
\gr

\medskip

\subsection*{Παρατηρήσεις}

\begin{itemize}
  \item Το σύστημα ακολουθεί με υψηλή ακρίβεια την επιθυμητή τροχιά $r_d(t)$ εντός του επιτρεπτού περιθωρίου $\pm\phi(t)$.
  \item Ο ελεγκτής έχει μικρές απαιτήσεις σε $u(t)$, χωρίς κορεσμό ή αστάθεια.
  \item Η λογαριθμική μορφή του ελεγκτή και το \en funnel \gr σύγκλισης εγγυώνται ομαλή και ασφαλή ρύθμιση.
  \item Η προσομοίωση είναι πολύ ακριβής.
\end{itemize}

\hspace{-0.6cm}Ο αντίστοιχος κώδικας \en MATLAB \gr βρίσκεται στο αρχείο \en Part2a.m\gr. 

\bigskip

\section{Θέμα 2(β) – Εκτίμηση παραμέτρων με μέθοδο \en Lyapunov \gr
         (μεικτή \en series–parallel \gr δομή)}

Το ίδιο μη‑γραμμικό μοντέλο
\[
\ddot{r}= -a_1\dot{r}-a_2\sin r
          +a_3\dot{r}^{2}\sin(2r)+b\,u
\]
με μετρούμενες καταστάσεις $(r,\dot r)$ και είσοδο $u(t)$  
υπόκειται στον ελεγκτή του Θέματος 2(α).  
Στόχος εδώ είναι η εκτίμηση σε πραγματικό χρόνο των άγνωστων 
\mbox{$a_1,a_2,a_3,b$} για $d(t)=0$.

\medskip

\subsection*{Δομή του εκτιμητή}

\paragraph{Παρατηρητής μεικτής μορφής}
\[
\begin{aligned}
\dot{\hat r}      &=\hat{\dot r}+ \theta_{M1}(r-\hat r) \\[2pt]
\dot{\hat{\dot r}}&=-\hat a_1\dot r-\hat a_2\sin r+\hat a_3\dot r^{2}\sin(2r)
                    +\hat b\,u+\theta_{M2}\bigl(\dot r-\hat{\dot r}\bigr)
\end{aligned}
\]
με $\theta_{M1}=50$, $\theta_{M2}=90$.

\paragraph{Νόμοι προσαρμογής}
\[
\begin{aligned}
\dot{\hat a}_1 &=-\gamma_1 \dot r\,(\dot r-\hat{\dot r}), &
\gamma_1=8.52 \\[-1pt]
\dot{\hat a}_2 &=-\gamma_2\sin r\,(\dot r-\hat{\dot r}), &
\gamma_2=0.0267 \\[-1pt]
\dot{\hat a}_3 &=\; \gamma_3\dot r^{2}\sin(2r)\,(\dot r-\hat{\dot r}), &
\gamma_3=0.01 \\[-1pt]
\dot{\hat b}   &=\;\gamma_4 u\,(\dot r-\hat{\dot r}), &
\gamma_4=0.03
\end{aligned}
\]

\hspace{-0.6cm}Η προσαρμογή βασίζεται σε κατάλληλη συνάρτηση \en Lyapunov \gr που 
περιλαμβάνει τα σφάλματα θέσης–ταχύτητας καθώς και τα σφάλματα παραμέτρων.  
Το σήμα
\(
\tilde{\dot r}=\dot r-\hat{\dot r}
\)
οδηγεί όλα τα $\hat a_i,\hat b$ να συγκλίνουν 
(\emph{τουλάχιστον} σε περιοχή) στις πραγματικές τιμές.

\medskip

\subsection*{Αποτελέσματα προσομοίωσης}

Ο κώδικας εκτελέστηκε για $t\in[0,20]\,\text{\en s\gr}$ με
βήμα $h=10^{-4}$\en s\gr.  
Τα τελικά αποτελέσματα:

\begin{center}
\begin{tabular}{@{}lcc@{}}
\toprule
Παράμετρος & Πραγματική & Τελική εκτίμηση \\
\midrule
$a_1$ & 1.3150 & 1.3168 \\
$a_2$ & 0.7250 & 0.7271 \\
$a_3$ & 0.2250 & 0.5001 \\
$b$   & 1.1750 & 1.1766 \\
\bottomrule
\end{tabular}
\end{center}

\begin{figure}[H]
  \centering
  \includegraphics[width=0.85\textwidth]{latex_and_pdf/2.b.2.png}
  \caption{Πραγματική και εκτιμώμενη γωνία $r(t)$, $\hat r(t)$}
\end{figure}

\begin{figure}[H]
  \centering
  \includegraphics[width=0.85\textwidth]{latex_and_pdf/2.b.3.png}
  \caption{Σφάλμα εκτίμησης $e_r(t)=r(t)-\hat r(t)$}
\end{figure}

\begin{figure}[H]
  \centering
  \includegraphics[width=0.85\textwidth]{latex_and_pdf/2.b.4.png}
  \caption{Εξέλιξη εκτιμήσεων $a_1(t)$, $a_2(t)$, $a_3(t)$, $b(t)$}
\end{figure}

\begin{figure}[H]
  \centering
  \includegraphics[width=0.85\textwidth]{latex_and_pdf/2.b.1.png}
  \caption{Σήμα ελέγχου $u(t)$ αναπαράγεται όπως Θέμα 2(α)}
\end{figure}

\medskip    

\subsection*{Παρατηρήσεις}

\begin{itemize}
  \item Τα $a_1$, $a_2$ και $b$ εκτιμώνται με σφάλμα πολύ μικρό μέσα σε
        $\approx 3$ \en sec\gr, επιβεβαιώνοντας την ταχεία σύγκλιση του εκτιμητή.
  \item Η μεγαλύτερη απόκλιση εμφανίζεται στο $a_3$ 
        ($\hat a_3 \simeq 0.50$) (αλλά με λίγο πιο σωστό \en tuning \gr των παραμέτρων προσαρμογής μου θα μπορούσε να συγκλινει στην επιθυμητή τιμή σωστά), 
        γεγονός αναμενόμενο γιατί ο αντίστοιχος όρος
        $\dot r^{2}\sin(2r)$ είναι μικρός στα πρώτα 5–6 \en sec \gr
        και δεν διεγείρει αρκετά τον προσαρμογέα.
  \item Το σφάλμα $e_r(t)$ παραμένει εντός του \en \textit{funnel} \gr
        που τέθηκε στο Θέμα 2(α) η ταυτόχρονη λειτουργία ελεγκτή 
        \& εκτιμητή δεν υποβαθμίζει την απόκριση.
  \item Η μεικτή δομή (\en series–parallel\gr) προσθέτει 
        τους όρους διόρθωσης $\theta_{M1},\theta_{M2}$,
        επιταχύνοντας την απόσβεση των σφαλμάτων
        χωρίς να απαιτεί μέτρηση επιτάχυνσης.
\end{itemize}

\hspace{-0.6cm}Ο αντίστοιχος κώδικας \en MATLAB \gr βρίσκεται στο αρχείο \en Part2b.m\gr.

\bigskip

\section{Θέμα 2(γ) – Εκτίμηση παραμέτρων υπό εξωτερική διαταραχή}

Στο Θέμα 2(γ) επαναλαμβάνεται η διαδικασία του ερωτήματος (β), εκτιμητής
\en Lyapunov \gr μεικτής δομής, με εξωτερική διαταραχή στην έξοδο:
\[
d(t)=0.15\,\sin(0.5\,t)\, .
\]

\medskip

\subsection*{Μοντέλο με διαταραχή}

Έτσι λοιπόν η εξίσωση μου δίνεται από την σχέση:
\[
\ddot{r}= -a_1\dot r-a_2\sin r+a_3\dot r^{2}\sin(2r)+b\,u+d(t)
\]
με $(a_1,a_2,a_3,b)=(1.315,\,0.725,\,0.225,\,1.175)$ όπως πριν.

\vspace{0.3cm}

\hspace{-0.6cm}Σε αντίθεση με τον εκτιμητή, ο οποίος αγνοεί το $d(t)$,
αυτό εισάγει σφάλμα μοντελοποίησης που επηρεάζει την ακρίβεια των
εκτιμήσεων.

\medskip

\subsection*{Ελεγκτής τροχιάς ίδιος με Θέμα 2(α), προσθήκη διαταραχής}

Ο ελεγκτής διατηρεί ακριβώς τον ίδιο μη‑γραμμικό νόμο με το
Θέμα 2(α).  Θέτουμε
\[
\phi(t)=\bigl(\phi_0-\phi_\infty\bigr)\,e^{-\lambda t}+\phi_\infty ,
\qquad
z_1(t)=\frac{r(t)-r_d(t)}{\phi(t)},
\quad
a(t)=-k_1\ln\!\Bigl(\frac{1+z_1(t)}{1-z_1(t)}\Bigr),
\]
\[
z_2(t)=\frac{\dot r(t)-a(t)}{\rho},
\qquad
u(t)=-k_2\ln\!\Bigl(\frac{1+z_2(t)}{1-z_2(t)}\Bigr),
\]
με
\(\phi_0=0.50\), \(\phi_\infty=0.01\), \(\lambda=0.5\),
\(\rho=1\), \(k_1=k_2=2\).
Το \emph{\en funnel \gr} $\phi(t)$ εγγυάται ότι
\(\lvert r(t)-r_d(t)\rvert<\phi(t)\)
παρά την ύπαρξη της διαταραχής
\(d(t)=0.15\sin(0.5\,t)\).

\medskip

\subsection*{Μεικτός εκτιμητής \en series–parallel \gr με \(d\ne0\)}

\paragraph{Παρατηρητής:}
Ο μεικτός παρατηρητής γράφεται
\[
\dot{\hat r}=\,\hat{\dot r}+\theta_{M1}\bigl(r-\hat r\bigr),
\]
\[
\dot{\hat{\dot r}}
     =-\hat a_1\dot r-\hat a_2\sin r+\hat a_3\dot r^{2}\sin(2r)
       +\hat b\,u
       +\theta_{M2}\bigl(\dot r-\hat{\dot r}\bigr),
\]
όπου \(\theta_{M1}=\theta_{M2}=10\).
Η πραγματική, άγνωστη δυναμική περιλαμβάνει τον επιπλέον όρο
\(d(t)\) συνεπώς το σφάλμα ταχύτητας
\(\tilde{\dot r}=\dot r-\hat{\dot r}\)
εξελίσσεται σύμφωνα με
\[
\dot{\tilde{\dot r}}
 =-a_1\dot r-a_2\sin r+a_3\dot r^{2}\sin(2r)+b\,u+d(t)
  +\theta_{M2}\tilde{\dot r}
  +\hat a_1\dot r+\hat a_2\sin r-\hat a_3\dot r^{2}\sin(2r)-\hat b\,u .
\]

\paragraph{Νόμοι προσαρμογής:}
Οι εκτιμήσεις ενημερώνονται με
\[
\dot{\hat a}_1=-\gamma_1\dot r\,\tilde{\dot r},\qquad
\dot{\hat a}_2=-\gamma_2\sin r\,\tilde{\dot r},\qquad
\dot{\hat a}_3=\; \gamma_3\dot r^{2}\sin(2r)\,\tilde{\dot r},\qquad
\dot{\hat b}\;=\; \gamma_4 u\,\tilde{\dot r},
\]
όπου επιλέχθηκαν
\[
\gamma_1=6,\quad
\gamma_2=0.01,\quad
\gamma_3=0.4,\quad
\gamma_4=0.03 .
\]

\hspace{-0.6cm}Οι παραπάνω σχέσεις προκύπτουν από συνάρτηση \en Lyapunov \gr
που περιλαμβάνει τα σφάλματα \(r-\hat r\), \(\dot r-\hat{\dot r}\)
και τα σφάλματα παραμέτρων επιβάλλοντας
\(\dot V\le 0\) προκύπτουν οι συγκεκριμένοι κανόνες προσαρμογής.

\medskip

\medskip

%----------------------------------------------------------------------
\subsection*{Αποτελέσματα προσομοίωσης ($h=10^{-4}\,$\en s\gr, $T=20\,$\en s\gr)}

\begin{center}
\begin{tabular}{@{}lcc@{}}
\toprule
Παράμετρος & Πραγματική & Τελική εκτίμηση \\ \midrule
$a_1$ & 1.3150 & 1.0317 \\
$a_2$ & 0.7250 & 0.7572 \\
$a_3$ & 0.2250 & 0.5101\,\\
$b$   & 1.1750 & 1.5700 \\ \bottomrule
\end{tabular}
\end{center}

\begin{figure}[H]
  \centering
  \includegraphics[width=0.85\textwidth]{latex_and_pdf/2.c.1.png}
  \caption{Έλεγχος: $u(t)$}
\end{figure}

\begin{figure}[H]
  \centering
  \includegraphics[width=0.85\textwidth]{latex_and_pdf/2.c.2.png}
  \caption{Πραγματική \& εκτιμώμενη γωνία $r(t)$, $\hat r(t)$ \;–\; $d(t)\neq0$}
\end{figure}

\begin{figure}[H]
  \centering
  \includegraphics[width=0.85\textwidth]{latex_and_pdf/2.c.3.png}
  \caption{Σφάλμα εκτίμησης $e_r(t)=r-\hat r$ (\textit{παραμένει φραγμένο})}
\end{figure}

\begin{figure}[H]
  \centering
  \includegraphics[width=0.85\textwidth]{latex_and_pdf/2.c.4.png}
  \caption{Εκτιμήσεις $a_1(t)$, $a_2(t)$, $a_3(t)$, $b(t)$ με διαταραχή}
\end{figure}

\medskip

\subsection*{Παρατηρήσεις}

\begin{itemize}
  \item Ο ελεγκτής διατηρεί το σφάλμα παρακολούθησης εντός του
        \en funnel \gr παρότι η έξοδος διεγείρεται από θόρυβο
        πλάτους $0.15\,\text{\en rad\gr}$.
  \item Έγινε προσπάθεια για εύρεση βέλιστων παραμέτρων γ αλλά δυστυχώς δεν        επιδιώχθηκε η σωστή σύγκλιση παρά την ορθότητα των εξισώσεων που           χρησιμοποιήθηκαν ενώ με κατάλληλο \en tuning \gr παραμέτρων θα             πετυχαίναμε την σύγκλιση στις αληθινές.      
\end{itemize}

\hspace{-0.6cm}Ο πλήρης κώδικας \en MATLAB \gr προσομοίωσης παρέχεται στο
\en Part2c.m \gr.



\bibliographystyle{plain}
\begin{thebibliography}{1}
    \bibitem{lnmpikas}
    \en https://elearning.auth.gr/course/view.php?id=16318
\end{thebibliography}

\end{document}